\documentclass[12pt,a4paper]{article}
\usepackage[utf8]{inputenc}
\usepackage[T1]{fontenc}
\usepackage{geometry}
\geometry{margin=2.5cm}
\usepackage{hyperref}
\usepackage{enumitem}
\usepackage{fancyhdr}
\usepackage{titlesec}
\usepackage{array}
\usepackage{longtable}
\usepackage{xcolor}

\title{\textbf{Spécifications des Interfaces à Développer - MineTrack}}
\author{Documentation Technique}
\date{\today}

\pagestyle{fancy}
\fancyhf{}
\rhead{MineTrack}
\lhead{Interfaces UI/UX}
\rfoot{Page \thepage}

\titleformat{\section}{\normalfont\Large\bfseries}{\thesection}{1em}{}[\titlerule]

\begin{document}


\section{Introduction}
Ce document détaille les interfaces utilisateur (UI) à développer pour le système \textbf{MineTrack}. Il définit les composants, les fonctionnalités et les interactions pour les plateformes Web et Mobile.

\section{Interfaces Application Web (Administration)}

\subsection{Tableau de Bord Global (Dashboard)}
\begin{itemize}[label=\textbullet]
    \item \textbf{Composants :}
    \begin{itemize}
        \item Cartes de statistiques (Production totale, Effectif présent, Alertes actives).
        \item Graphique linéaire (Évolution de la production hebdomadaire).
        \item Liste des derniers incidents signalés avec indicateurs de couleur.
    \end{itemize}
    \item \textbf{Interactions :} Filtrage par site minier, survol des graphiques pour détails.
\end{itemize}

\subsection{Gestion du Personnel et Plannings}
\begin{itemize}[label=\textbullet]
    \item \textbf{Composants :}
    \begin{itemize}
        \item Tableau de données (Nom, Rôle, Site, Statut).
        \item Calendrier de shift interactif (Grille semaine/mois).
        \item Formulaire de création/édition en modal.
    \end{itemize}
    \item \textbf{Fonctionnalités :} Recherche par matricule, assignation rapide de shift par drag-and-drop.
\end{subsection}

\subsection{Module de Rapports (Production \& Sécurité)}
\begin{itemize}[label=\textbullet]
    \item \textbf{Composants :}
    \begin{itemize}
        \item Filtres multicritères (Date, Site, Type, Gravité).
        \item Visionneuse de rapports détaillée.
        \item Boutons d'exportation (PDF, Excel).
    \end{itemize}
\end{itemize}

\section{Interfaces Application Mobile (Terrain)}

\subsection{Écran de Pointage (Chef d'Équipe)}
\begin{itemize}[label=\textbullet]
    \item \textbf{Composants :}
    \begin{itemize}
        \item Liste défilante des ouvriers assignés au shift.
        \item Boutons d'action rapides (Check vert pour présent, Croix rouge pour absent, Horloge pour retard).
        \item Badge indicateur pour les employés en repos (OFF).
    \end{itemize}
    \item \textbf{Interactions :} Recherche par nom, validation globale du shift.
\end{itemize}

\subsection{Saisie de Production}
\begin{itemize}[label=\textbullet]
    \item \textbf{Composants :}
    \begin{itemize}
        \item Sélecteur de zone (Boutons segmentés).
        \item Champs de saisie numérique (Tonnage, Nombre de camions).
        \item Bouton de soumission haute visibilité.
    \end{itemize}
\end{itemize}

\subsection{Déclaration d'Incident}
\begin{itemize}[label=\textbullet]
    \item \textbf{Composants :}
    \begin{itemize}
        \item Menu déroulant (Type d'incident).
        \item Sélecteur de gravité visuel (Échelle colorée de Vert à Rouge).
        \item Zone de texte pour description et localisation précise.
    \end{itemize}
\end{itemize}

\subsection{Planning Personnel (Ouvrier)}
\begin{itemize}[label=\textbullet]
    \item \textbf{Composants :}
    \begin{itemize}
        \item Vue hebdomadaire simplifiée.
        \item Détails du shift (Heure, Zone, Équipement assigné).
        \item Système de notifications intégré.
    \end{itemize}
\end{itemize}

\section{Charte Graphique \& UX}
\begin{itemize}[label=\textbullet]
    \item \textbf{Couleurs de Gravité :}
    \begin{itemize}
        \item \textcolor{green}{Faible (Vert)}
        \item \textcolor{yellow}{Moyen (Jaune)}
        \item \textcolor{orange}{Élevé (Orange)}
        \item \textcolor{red}{Critique (Rouge)}
    \end{itemize}
    \item \textbf{Mode Sombre/Clair :} Support natif sur les deux plateformes.
    \item \textbf{Accessibilité :} Taille de police ajustable et contrastes élevés pour usage en extérieur.
\end{itemize}

\end{document}
